\chapter{Masques et retouche locale}
L’onglet \emph{Masques}, en haut du panneau de droite, applique des réglages à une partie de
l’image seulement. Un masque est une zone qui porte ses propres valeurs : elles s’ajoutent aux
réglages globaux dans cette zone. Une exposition de +1 dans un masque agit comme une exposition
globale de +1, mais là seulement. Rien n’est écrit dans la photo : les masques s’enregistrent dans
son compagnon et s’annulent comme les autres étapes.

Les masques n’ont ni ordre ni mode de fusion : où plusieurs se recouvrent, leurs valeurs
s’additionnent.

\section{Créer un masque}
\begin{itemize}
\item \emph{Pinceau} : peignez sur l’image. \emph{Taille}, \emph{Contour progressif} et
\emph{Débit} règlent le coup de pinceau suivant ; chaque coup est une étape de l’historique.
\item \emph{Dégradé linéaire} : l’effet est entier avant la première poignée et nul après la seconde.
\item \emph{Dégradé radial} : déplacez le centre et les deux bords sur l’image. \emph{Contour
progressif} règle la transition, \emph{Angle} tourne l’ellipse.
\item \emph{Plage} : choisit des tons de l’image, entre \emph{Tons sombres} et \emph{Tons clairs},
tels qu’ils sont avant toute retouche. Une nouvelle plage couvre toute l’image : resserrez ces
curseurs pour isoler des tons.
\end{itemize}
\emph{Afficher la zone} montre où le masque agit : en couleur dedans, en gris dehors. Cet affichage
n’est ni enregistré ni exporté. La zone s’affiche aussi tant que vous tenez un curseur de la plage.
\emph{Inverser} retourne une forme.

\section{Composer une zone}
\emph{Ajouter à la zone} combine une autre forme à la première : \emph{Ajouter}, \emph{Soustraire}
ou \emph{Intersecter}. Une plage intersectée avec un dégradé donne, par exemple, les ombres du
seul bas de l’image.

\section{Régler un masque}
Chaque masque propose les réglages de lumière, de couleur, d’effets et de détail, ainsi qu’une
\emph{Opacité} qui pondère le tout. La case de la liste le désactive sans le perdre, et les deux
boutons voisins le dupliquent ou le suppriment. Un masque suit la géométrie de l’image : il reste
sur le même objet après une rotation ou un recadrage.

Les masques appartiennent à leur photo : coller des réglages ou les appliquer à une sélection ne
les transmet pas. \emph{Réinitialiser les réglages} les retire aussi.

\section{Retirer un objet}
L’onglet \emph{Retouche} retire ce qui ne doit pas être dans la photo. Ce qui remplace la zone est
reconstruit, pas retrouvé. La photo d’origine ne change pas : chaque retouche s’enregistre dans son
compagnon, peut être masquée, modifiée ou retirée, et s’annule comme les autres étapes.

Choisissez une action, réglez la taille du pinceau, puis peignez sur l’image :
\begin{itemize}
\item \emph{Effacer} invente la zone à partir de ce qui l’entoure. C’est convaincant sur un ciel ou
une texture régulière ; regardez le résultat à 100~\% ailleurs.
\item \emph{Corriger} prend la texture d’un autre endroit, à la lumière des alentours de la zone.
C’est l’outil d’une poussière ou d’un petit défaut.
\item \emph{Tampon} copie un autre endroit tel quel.
\end{itemize}
Pour \emph{Corriger} et \emph{Tampon}, une source est proposée au premier coup de pinceau ;
\emph{Déplacer la source} montre une poignée pour en choisir une autre. Si une zone est trop grande
pour être reconstruite, ou si sa source sort de l’image, l’image reste telle quelle et le message
l’indique.

Faites le redressement d’abord : comme celle d’un masque, la zone ne suit pas l’image si la géométrie
change ensuite.

\section{Limites}
\begin{itemize}
\item Le bruit de luminance d’un masque est traité après le profil, même en développement signé.
\item Modifier un masque est plus lent quand il y en a beaucoup : au-delà d’une quinzaine, le geste
dépasse le quart de seconde. Bouger un autre curseur ne les redessine pas.
\end{itemize}
